本文以社区非隔离感染人数为直接控制对象，要求全过程满足 $I(t)\le\eta$，其中 $\eta>0$ 为给定上限，并记 $\theta=\eta/N$。Miclo 等\cite{Miclo2022ICU}通过比例关系将感染人数上限与 ICU 容量联系起来；Angulo 等\cite{Angulo2021}则先给定允许现患率，再在应用中依据医疗资源说明其取值。由于本文区分非隔离与隔离感染者，容量解释还需要计入两类病例的共同来源。以下说明社区感染上限与医疗容量的联系，再研究给定阈值下的控制策略。

\begin{remark}
在模型\eqref{eq:model:full}中，两类感染仓室的新增流入之和为
\[
\frac{\beta c(t)(1-q(t))S(t)I(t)}{N}
+\frac{\beta c(t)q(t)S(t)I(t)}{N}
=\frac{\beta c(t)S(t)I(t)}{N}.
\]
若全过程有 $0\le c(t)\le c_0$ 和 $I(t)\le\eta$，则由 $S(t)\le S_0$ 得
\[
0\le\frac{\beta c(t)S(t)I(t)}{N}
\le\frac{\beta c_0S_0}{N}\eta.
\]
这给出全部新增感染流量的统一上界，而不只约束留在社区的病例。

为说明该上界与医疗容量的联系，附加如下医疗需求关系。设 $h(a)\ge0$ 为一例新感染者在感染后第 $a$ 天所需的平均 ICU 床位数；假定该函数可测，且对进入 $I$、进入 $I_q$ 的两类病例、不同感染时点及所比较的干预分配共同适用。设
\[
\int_0^\infty h(a)\,da=pL<\infty,
\]
其中 $p\in(0,1]$ 是一例模型新感染者此后需要 ICU 的概率，$L>0$ 是需要 ICU 者累计占用时长的均值。$h$ 可以包含感染至入院的延迟，不要求指定其具体分布。由既往流入与停留过程计算床位占用的建模方法可参见文献\cite{Baas2021Occupancy}。

记初始病例在此后产生的平均需求为 $D_0(t)$，并假定对所比较策略均有 $0\le D_0(t)\le D_{\mathrm{init}}<\infty$。该项包括初始时刻已经感染但以后才需要 ICU 的病例，不能仅用初始在床人数代替。未受容量拒收限制的平均需求定义为
\[
D(t)=D_0(t)+\int_0^t h(t-u)
\frac{\beta c(u)S(u)I(u)}{N}\,du.
\]
由非负性及前述流量界，直接得到
\[
\begin{aligned}
D(t)
&\le D_{\mathrm{init}}+
\frac{\beta c_0S_0}{N}\eta\int_0^t h(a)\,da\\
&\le D_{\mathrm{init}}+pL\frac{\beta c_0S_0}{N}\eta,
\qquad t\ge0.
\end{aligned}
\]
\Needspace{9\baselineskip}
设 $C>D_{\mathrm{init}}$ 为扣除其他病种占用和安全预留后可用于本病的容量，其中本病初始病例的负担由 $D_{\mathrm{init}}$ 单独计入、不重复扣除。因而
\[
\eta\le\frac{C-D_{\mathrm{init}}}{pL\,\beta c_0S_0/N}
\]
是上述平均需求满足 $D(t)\le C$ 的充分条件。两类病例均已计入总流量，不需要假定隔离感染者没有重症风险。该条件覆盖控制前、控制期和退出后；它不要求需求峰值与感染峰值同时出现，也不表示 $I=\eta$ 时 ICU 恰好满载。这里的结论针对所规定的平均需求关系，不是实际随机占用的逐样本保证；不满足该充分条件也不能据此判定必然超容。
\end{remark}

上述医疗需求关系用于解释阈值的选取，不加入传播方程、人口守恒式或成本目标。本文不估计 $h,p,L$ 或初始病例的医疗需求；应用中的比例参考另行说明，不将其与这里的充分条件等同。以下对满足初值、触发及实施条件的给定 $\eta$ 分析三阶段控制。各策略使用同一阈值，不按各自的感染分流结果重新调整。
